\chapter{Sources}

Each source is listed with the claims this book takes from it. No link is given, because a link is a
statement that the source was retrieved and read, and the honest state is that these are the
foundational works named from knowledge of the field and not yet opened, entered and identified by
\texttt{maint/citations/citations.py}. An empty identifier means it has not been looked up yet, which
is the same discipline the provenance page states for the rest of this work.

\section*{The order and the periodic law}

\textbf{D. Mendeleev}, the periodic law, 1869. A property of the element is a periodic function of its
place in the order. This book takes the periodic reading of a property along $Z$ from it.

\textbf{H. Moseley}, the high-frequency spectra of the elements, 1913. The order is the nuclear
charge $Z$ and not the atomic weight. This book takes from it that the ordinate the property is read
against is a count of protons.

\section*{The shell counts}

\textbf{W. Pauli}, the exclusion principle, 1925. No two electrons share all quantum numbers, which
is why a shell $n$ holds $2n^2$ states. This book takes the shell capacities $2, 8, 18, 32$ and the
period lengths that follow from them as fixed before any measurement.

\section*{Valence and the octet}

\textbf{G. N. Lewis}, the atom and the molecule, 1916. The shared electron pair and the tendency of
an atom toward a completed outer shell. This book takes the octet as a condition on every atom of an
arrangement.

\textbf{I. Langmuir}, the arrangement of electrons in atoms and molecules, 1919. The octet named, and
covalence counted as shared pairs. This book takes the boolean reading at each atom from it.

\section*{Electronegativity and the bond}

\textbf{L. Pauling}, the nature of the chemical bond, 1932 and the later book. The electronegativity
scale and the tabulation of bond lengths and covalent radii. This book takes the element vector's
electronegativity component and the fact that a bond length is fixed by the bond type from it.

\section*{The value tables, not yet entered}

The specific magnitudes this book quotes are standard reference values and each needs its own
verified entry: the covalent and van der Waals radii, the Pauling electronegativities, and the bond
lengths, including the backbone values near $1.46$, $1.52$ and $1.33$ angstroms the protein reader
already consults. They are held in crystallographic and gas-phase spectroscopic tabulations and in
the IUPAC recommendations, and the row for each stays empty here until it is opened and identified.
Naming a value without the table it came from is the defect the citation tool exists to catch, and
this section states the debt and does not hide it.
